% SGA 3, 第 VI_B 篇, 第 1 轮英文重构底本的中译.
% 范围: 引理 7.5 及其完整证明, 包括原编者注 75.
% 权威底本: Polo--Gille Exp6B-13oct24.pdf, 组件第 56 页
% (合卷第 397 页; 重版页眉第 386 页; 原版边栏页码第 370 页).
% 注 7.6.0 始于组件第 56 页稍后位置, 不在本组件内.
% OCR 与 Reinhold 英文 Markdown 仅为对照材料. 本组件无图.

\medskip
\phantomsection\label{lem:sga3-VIB-7.5}%
\noindent\textbf{引理 7.5.---}
\emph{设 \(S\) 是一个概形, \(G\) 是一个 \(S\)-群概形,
\(X\) 是一个 \(S\)-预概形, \(f:X\to G\) 是一个 \(S\)-态射.
假设 \(X\) 与 \(G\) 在 \(S\) 上局部有限呈示,
且对每个 \(s\in S\) 及 \(G_s\) 的每个极大点 \(g\),
存在 \(X\) 的一个点 \(x\), 使得 \(f(x)=g\),
且 \(f\) 在 \(x\) 处平坦. 则态射}
\phantomsection\label{form:sga3-VIB-7.5-covering-morphism}%
\[
  \mu\circ(f\times_S f):
  X\times_S X\longrightarrow G
\]
\emph{是 \((\mathrm{fppf})\) 拓扑的一个覆盖.}

\footnote{\label{note:sga3-VIB-75}编者注: 以下扩充了原文的论证.}
由 EGA~IV\(_3\), 11.3.1, \(X\) 中 \(f\) 在其处平坦的那些点
所成的集合 \(V\) 是开的, 且 \(f|_V\) 是开态射. 所以
\[
  U=f(V)
\]
是 \(G\) 的一个开子集. 此外, 由假设, 对每个 \(s\in S\),
\[
  U\cap G_s
\]
在 \(G_s\) 中稠密.
以 \(\phi\) 表示 \(\mu\circ(f\times_S f)\)
在 \(V\times_S V\) 上的限制.

只需证明 \(\phi\) 是 \((\mathrm{fppf})\) 拓扑的一个覆盖;
为此只需证明 \(\phi\) 忠实平坦且有限呈示
(\hyperref[sga3:iv:proposition:6:3:1:item:iv:01]
{第 IV 篇, 6.3.1\textup{(iv)}}).
现在, \(\phi\) 是复合
\phantomsection\label{form:sga3-VIB-7.5-factorization}%
\[
  V\times_S V
  \xrightarrow{\ f|_V\times f|_V\ }
  U\times_S U
  \xrightarrow{\ \mu\ }
  G,
\]
其中第一个态射忠实平坦且局部有限呈示, 因为 \(f|_V\) 具有这些性质.
所以只需再对 \(\mu\) 在 \(U\times_S U\) 上的限制证明相同的断言.

由于 \(G\to S\) 局部有限呈示且平坦, 态射
\[
  \mu:G\times_S G\longrightarrow G.
\]
也如此. 事实上, 此态射同构于将 \(G\to S\)
沿基变换 \(G\to S\) 拉回所得的态射
(\hyperref[sga3:VIA:0.1]{VI\(_A\), 0.1}).
因此, 诱导的态射
\[
  U\times_S U\longrightarrow G
\]
有相同的性质. 要证明最后这个态射为满射, 只需逐纤维核验.
在每个纤维上, 结果由 \hyperref[sga3:VIA:0.5]{VI\(_A\), 0.5} 得到,
因为对每个 \(s\in S\), \(U\cap G_s\) 是 \(G_s\) 的一个稠密开子集.

% 中文翻译检查点：本组件全部正文、证明、脚注及源注已译；源79行。
% 机械核对：行内数学45处、陈列数学6处及全部引用键守恒（数学中的自然语言已译）。
% 状态：完成源文覆盖和逐段自校；尚未经独立合卷编译与页面审查。
